\documentclass[11pt]{article}
\usepackage[margin=1in]{geometry}
\usepackage{amsmath}
\title{Four Pillars of Organizational Resilience: A Conceptual Framework}
\author{Alicia Reyes \and Hamid Al-Qasemi}
\date{2023}
\begin{document}
\maketitle
\noindent DOI (synthetic): 10.9999/trap.2024.005

\begin{abstract}
We propose a four-pillar framework for organizational resilience: (1) adaptive capacity, (2) redundancy, (3) information sharing, and (4) leadership continuity. Each pillar is conceptually distinct; their combination is theorized to be necessary and sufficient for resilience.
\end{abstract}

\section{Introduction}
Resilience is widely discussed but poorly structured. We synthesize the prior literature into four orthogonal constructs.

\section{The Four Pillars}

\subsection{Adaptive Capacity}
Adaptive capacity is the organization's ability to reconfigure routines and resource allocations in response to shocks. Distinguished from learning (which is retrospective) by its prospective, real-time reconfiguration focus.

\subsection{Redundancy}
Redundancy is the presence of slack resources and parallel capabilities that permit continued operation when individual components fail. Distinguished from inefficiency: redundancy is purposeful duplication; inefficiency is accidental.

\subsection{Information Sharing}
Information sharing is the organization-wide practice of rapidly disseminating weak signals about threats, opportunities, and deviations. Distinguished from communication: information sharing emphasizes signal speed and breadth over structure.

\subsection{Leadership Continuity}
Leadership continuity is the preservation of decision authority through turnover, absence, or disruption. Distinguished from hierarchy: continuity may be achieved via flat succession arrangements.

\section{Framework Proposition}
We propose that organizational resilience $R$ is a function of all four pillars jointly:
\[
R = f(\text{AC}, \text{Red}, \text{IS}, \text{LC})
\]
Each pillar is necessary; no three pillars alone are sufficient. Removal of any pillar collapses resilience in the face of sufficiently large shocks.

\section{Discussion}
Empirical operationalization of each pillar requires further work. We offer this framework as a scaffold for that effort.

\end{document}
